\documentclass[12pt,a4paper]{article}
\usepackage[UTF8]{ctex}
\usepackage{geometry}
\usepackage{booktabs}
\usepackage{amsmath}
\usepackage{setspace}
\usepackage{hyperref}
\usepackage{xcolor}
\geometry{left=2.7cm,right=2.7cm,top=2.5cm,bottom=2.5cm}
\onehalfspacing
\hypersetup{colorlinks=true,linkcolor=blue,urlcolor=blue,citecolor=blue}

\title{审计师行业专长与企业漂绿：两轮回归与关键诊断报告}
\author{}
\date{\today}

\begin{document}
\maketitle
\tableofcontents
\newpage

\section{研究目的与结论摘要}

本报告汇报审计师行业专长与企业漂绿关系的两轮实证分析。研究使用2021--2024年中国上市公司面板数据，主要考察会计师事务所在特定行业积累的审计经验能否抑制客户企业漂绿。第一轮采用事务所当期行业客户资产市场份额衡量行业专长；第二轮针对该指标可能包含客户企业自身资产而产生机械相关的问题，构造剔除本企业后的市场份额（leave-one-out，以下简称LOO），并进一步使用客户数量份额和公司固定效应模型进行诊断。

两轮结果表明，原始资产市场份额与漂绿程度呈负向且边际显著的关系，但剔除企业自身资产后，该关系几乎完全消失；LOO客户数量份额保持负向，但统计上不显著；加入公司固定效应后，两种LOO指标均接近于零。因此，当前最审慎的阶段性判断是：

\begin{quote}
现有结果尚未提供审计师行业专长显著抑制企业漂绿的稳健证据。原始资产市场份额的负向结果，可能主要受到企业自身规模机械进入专长指标的影响，而不宜直接解释为审计师行业知识的监督效应。
\end{quote}

上述结论并不意味着研究假设必然不成立，而是说明当期横截面资产市场份额可能不是最合适的专长度量。下一阶段应优先考察历史累计行业经验、滞后专长指标和事务所变更事件。

\section{数据、样本与变量}

\subsection{样本与数据处理}

原始面板包含9,370个公司--年度观测，覆盖2,915家公司，样本期为2021--2024年。数据合并了企业漂绿、财务指标、审计机构、审计意见、股权结构、公司治理和上市日期等信息。

主要数据处理包括：
\begin{itemize}
  \item 以六位证券代码和年度识别公司--年度观测；
  \item 对连续漂绿指标及连续控制变量进行上下1\%缩尾；
  \item 企业规模为总资产自然对数；资产负债率为总负债除以总资产；
  \item 总资产收益率为净利润除以总资产；经营现金流为经营活动现金流除以总资产；
  \item 第一大股东持股比例按百分数转换为比例；
  \item 企业年龄为$\ln(1+\text{样本年度}-\text{上市年份})$；
  \item 标准误均在公司层面聚类。
\end{itemize}

连续漂绿指标缺失79条。完整控制变量样本包含8,577个观测和2,746家公司，样本减少主要来自实际控制人性质和两职合一变量缺失。第二轮固定样本比较统一使用这8,577条观测，以避免样本变化干扰模型之间的系数比较。

\subsection{被解释变量与控制变量}

主要被解释变量为连续漂绿指标\texttt{Greenwashing连续\_基准}。第一轮还使用二元漂绿变量进行替代检验。

控制变量包括企业规模、资产负债率、总资产收益率、经营现金流、第一大股东持股比例、国有企业虚拟变量、董事长与总经理两职合一、独立董事比例、企业年龄以及国际四大审计虚拟变量。

\section{审计师行业专长的测量}

\subsection{原始资产市场份额}

第一轮使用事务所在行业年度内审计客户的资产市场份额：
\begin{equation}
AssetShare_{ajt}=
\frac{\sum_{i\in a,j}Assets_{it}}
{\sum_{i\in j}Assets_{it}},
\end{equation}
其中$a$表示会计师事务所，$j$表示证监会行业，$t$表示年度。第一轮还基于该份额构造10\%行业专家虚拟变量和行业市场份额第一名虚拟变量。

该指标的潜在问题是，企业$i$自身资产同时进入分子和分母。当某家公司规模较大时，即使事务所在该行业缺少其他客户，该企业本身也可能使其审计师获得较高市场份额。因此，原始指标可能同时反映事务所行业经验和客户自身规模。

\subsection{LOO资产市场份额}

为排除企业自身造成的机械影响，第二轮构造：
\begin{equation}
AssetShare^{LOO}_{iajt}=
\frac{AuditorIndustryAssets_{ajt}-Assets_{it}}
{IndustryAssets_{jt}-Assets_{it}}.
\end{equation}

该指标衡量：不考虑企业自身时，其会计师事务所在相同行业拥有的其他客户资产份额。相较原始资产份额，LOO指标更接近事务所从其他客户中积累的行业经验。

\subsection{LOO客户数量份额}

为降低少数大客户对资产口径的影响，进一步构造：
\begin{equation}
ClientShare^{LOO}_{iajt}=
\frac{AuditorIndustryClients_{ajt}-1}
{IndustryClients_{jt}-1}.
\end{equation}

该指标衡量事务所在行业内审计其他客户的数量占比，更侧重审计经验的广度。

\section{计量模型与诊断设计}

行业和年度固定效应模型为：
\begin{equation}
Greenwashing_{it}=\beta Expertise_{it}+\gamma Controls_{it}
+\delta_j+\lambda_t+\varepsilon_{it},
\end{equation}
其中$\delta_j$为行业固定效应，$\lambda_t$为年度固定效应。

第二轮进一步估计公司固定效应模型：
\begin{equation}
Greenwashing_{it}=\beta Expertise^{LOO}_{it}+\gamma Controls_{it}
+\mu_i+\lambda_t+\varepsilon_{it},
\end{equation}
其中$\mu_i$为公司固定效应。该模型利用同一家企业随时间的专长变化进行识别，并吸收企业长期不变的文化、管理质量、环境战略和信息透明度等不可观测特征。

第二轮诊断主要回答三个问题：第一，第一轮系数变化是否只是样本变化造成；第二，原始资产份额结果是否来自企业自身规模的机械影响；第三，利用公司内部变化后是否仍能识别行业专长的作用。

\section{第一轮回归结果}

\begin{table}[htbp]
\centering
\caption{第一轮核心回归结果}
\label{tab:first}
\small
\begin{tabular}{lccccc}
\toprule
 & M1 & M2 & M3 & M4 & M5 \\
\midrule
行业专家（10\%） & $-0.074^{*}$ & $-0.013$ & & & $0.003$ \\
 & $(0.042)$ & $(0.044)$ & & & $(0.013)$ \\
原始资产市场份额 & & & $-0.473^{*}$ & & \\
 & & & $(0.244)$ & & \\
行业排名第一 & & & & $0.001$ & \\
 & & & & $(0.061)$ & \\
控制变量 & 否 & 是 & 是 & 是 & 是 \\
年度固定效应 & 是 & 是 & 是 & 是 & 是 \\
行业固定效应 & 是 & 是 & 是 & 是 & 是 \\
观测值 & 9,291 & 8,577 & 8,577 & 8,577 & 8,647 \\
调整$R^2$ & $-0.007$ & 0.007 & 0.008 & 0.007 & 0.037 \\
\bottomrule
\end{tabular}

\begin{minipage}{0.93\textwidth}
\footnotesize 注：括号内为公司层面聚类标准误。$^{*}$、$^{**}$和$^{***}$分别表示在10\%、5\%和1\%水平显著。M5以二元漂绿为被解释变量，采用线性概率模型。
\end{minipage}
\end{table}

M1中10\%行业专家虚拟变量为负且在10\%水平显著；加入控制变量后，M2的系数下降至$-0.013$且不显著。M3中原始资产市场份额系数为$-0.473$，$p=0.053$，是第一轮最接近支持研究假设的结果。行业第一指标和二元漂绿模型均不显著。这表明第一轮证据较弱，并且依赖连续资产市场份额这一特定度量。

\section{第二轮关键诊断结果}

\subsection{固定样本逐步回归}

\begin{table}[htbp]
\centering
\caption{固定样本下原始资产市场份额的逐步回归}
\label{tab:same-sample}
\small
\begin{tabular}{lccc}
\toprule
 & 无公司控制 & 财务控制 & 完整控制 \\
\midrule
原始资产市场份额 & $-0.896^{***}$ & $-0.494^{**}$ & $-0.473^{*}$ \\
 & $(0.237)$ & $(0.244)$ & $(0.244)$ \\
$p$值 & 0.0002 & 0.0427 & 0.0530 \\
观测值 & 8,577 & 8,577 & 8,577 \\
年度固定效应 & 是 & 是 & 是 \\
行业固定效应 & 是 & 是 & 是 \\
\bottomrule
\end{tabular}

\begin{minipage}{0.93\textwidth}
\footnotesize 注：所有模型使用相同的8,577条完整样本。财务控制包括规模、资产负债率、总资产收益率和经营现金流；完整控制进一步加入股权、治理、年龄及国际四大变量。
\end{minipage}
\end{table}

固定样本结果表明，第一轮系数下降并非主要由样本变化造成。加入财务控制变量后，系数绝对值由0.896下降到0.494，说明企业财务特征解释了原始负相关关系中的较大部分；继续加入股权和治理变量后，系数变化较小。该现象提示企业规模等财务特征可能同时影响漂绿与事务所资产市场份额。

\subsection{LOO指标与公司固定效应}

\begin{table}[htbp]
\centering
\caption{剔除企业自身后的关键诊断回归}
\label{tab:loo}
\small
\begin{tabular}{lcccccc}
\toprule
 & L1 & L2 & L3 & C1 & F1 & F2 \\
\midrule
LOO资产份额 & $-0.086$ & $0.001$ & $-0.010$ & & $0.012$ & \\
 & $(0.264)$ & $(0.268)$ & $(0.269)$ & & $(0.277)$ & \\
LOO客户份额 & & & & $-0.330$ & & $-0.033$ \\
 & & & & $(0.331)$ & & $(0.416)$ \\
$p$值 & 0.746 & 0.996 & 0.971 & 0.319 & 0.966 & 0.937 \\
财务控制 & 否 & 是 & 是 & 是 & 是 & 是 \\
治理等控制 & 否 & 否 & 是 & 是 & 是 & 是 \\
公司固定效应 & 否 & 否 & 否 & 否 & 是 & 是 \\
行业固定效应 & 是 & 是 & 是 & 是 & 否 & 否 \\
年度固定效应 & 是 & 是 & 是 & 是 & 是 & 是 \\
观测值 & 8,577 & 8,577 & 8,577 & 8,577 & 8,577 & 8,577 \\
\bottomrule
\end{tabular}

\begin{minipage}{0.96\textwidth}
\footnotesize 注：括号内为公司层面聚类标准误。L1--L3依次加入控制变量；C1使用LOO客户数量份额；F1和F2加入公司与年度固定效应。公司固定效应模型不另行加入行业固定效应，因为公司固定效应已吸收企业不随时间变化的行业归属。
\end{minipage}
\end{table}

最关键的变化是：原始资产市场份额在完整控制模型中的系数为$-0.473$，而LOO资产市场份额在相同样本和控制变量下的系数仅为$-0.010$，$p=0.971$。系数绝对值接近于零，说明剔除企业自身资产后，原始负向关系基本消失。

LOO客户数量份额在行业和年度固定效应模型中的系数为$-0.330$，方向符合预期，但$p=0.319$，无法拒绝系数为零的假设。公司固定效应模型中，LOO资产份额和LOO客户份额的系数分别为0.012和$-0.033$，均接近于零且完全不显著。

\subsection{公司内变化与事务所变更}

\begin{table}[htbp]
\centering
\caption{行业专长指标的公司内变化}
\label{tab:variation}
\small
\begin{tabular}{lccc}
\toprule
指标 & 总体标准差 & 公司内标准差 & 指标发生变化的公司数 \\
\midrule
原始资产份额 & 0.108 & 0.045 & 2,477 \\
LOO资产份额 & 0.088 & 0.047 & 2,189 \\
LOO客户数量份额 & 0.072 & 0.032 & 2,168 \\
\bottomrule
\end{tabular}
\end{table}

完整样本中共有2,746家公司，其中734家公司发生过事务所更换，共识别818次变更事件。LOO指标在超过2,100家公司中存在至少两个不同取值，因此，公司固定效应不显著不能完全归因于解释变量没有时间变化。不过，由于样本期仅四年，公司内识别空间仍然有限。

原始资产份额与LOO资产份额的相关系数为0.757，LOO资产份额与LOO客户数量份额的相关系数为0.775。两种LOO口径具有较高但并非完全一致的相关性，说明资产规模和客户经验广度分别包含不同信息。

\section{综合判断}

第二轮诊断改变了对第一轮结果的解释。第一轮原始资产市场份额的负向边际显著结果不能稳健地归因于审计师行业知识，原因如下：

\begin{enumerate}
  \item 在固定样本中加入企业财务特征后，原始份额系数明显缩小，说明客户财务特征解释了较大部分相关关系；
  \item 剔除企业自身资产后，LOO资产份额的系数由$-0.473$下降到$-0.010$，且$p$值升至0.971；
  \item LOO客户数量份额虽保持负向，但统计上不显著；
  \item 控制公司不随时间变化的不可观测特征后，两种LOO指标均接近于零；
  \item 第一轮10\%阈值、行业第一和二元漂绿模型本身也未形成一致证据。
\end{enumerate}

因此，目前不能将研究结论表述为“审计师行业专长显著抑制企业漂绿”，也不宜仅保留原始资产市场份额的边际显著结果。较为规范的阶段性表述是：

\begin{quote}
传统的当期客户资产市场份额与企业漂绿程度呈负相关，但该关系在剔除客户企业自身后消失，表明原始结果可能受到客户规模造成的机械关联影响。基于更严格的LOO指标和公司固定效应，尚未发现审计师行业专长显著影响企业漂绿的证据。
\end{quote}

\section{阶段性结论}

两轮分析的主要价值在于识别出原始行业专长度量中的机械关联问题。第一轮结果显示原始资产市场份额与漂绿程度负向相关，但第二轮LOO诊断表明，该关系主要由客户企业自身资产驱动。当前更严格的结果不支持审计师行业专长显著抑制企业漂绿。因此，下一阶段不应围绕原始边际显著结果进行结论包装，而应改进专长指标，使其更准确地反映事务所从其他客户和历史业务中积累的行业知识，再判断研究假设是否成立。

\end{document}
